% Ch3 self-study (part 1) -- "I do / You do" ladder for Zumdahl 3.1-3.6.
% Companion to ch03-selfstudy (3.7-3.11). Same format: worked example,
% then a twin problem whose bare final answer is visible in every edition
% (\selfcheck) while the full solution stays key-only. Examples disjoint
% from ch03-notes and both other Ch3 documents.
\documentclass{apchem-worksheet}

\apcunitword{Chapter}
\apcunit{3}
\apcunittitle{Stoichiometry}
\apcdoctitle{Self-Study \textbullet\ \S3.1--3.6, I~do / You~do}
\apcsubtitle{Zumdahl \S3.1--3.6 \textbullet\ PDF pp.\ 111--124 \textbullet\ work every YOUR TURN before moving on}

\begin{document}
\apctitleblock

\begin{apcnote}[How to use these notes --- read this first]
Each skill comes as a \textbf{ladder} with two rungs. The solid-framed box
is the \textbf{worked example}: read it slowly, with a pencil, reproducing
each step. The dashed box is \textbf{YOUR TURN}: the same problem with new
numbers, and no help.

\begin{enumerate}[leftmargin=*,itemsep=0.15em]
  \item \textbf{Attempt the YOUR TURN completely} before looking at
        anything. Write real work in the workspace, not fragments.
  \item Check your result against the small gray \emph{check:} line. Right
        first try? Tick it on the tracker (last page) and climb on.
  \item Wrong? \textbf{Re-read the worked example, then redo} --- do not
        stare at your first attempt. Full solutions are in the teacher key
        if you are stuck after two tries.
\end{enumerate}
These six sections build one idea: the \textbf{mole} is how chemists count
particles by weighing --- and everything else in Chapter 3 stands on it.
\end{apcnote}

% =====================================================================
\section*{\sffamily Ladder 1 \textbullet\ Counting by weighing; average
atomic mass}
\zsec{3.1--3.2}

You cannot count atoms one at a time, but you can \emph{weigh} them in
bulk --- and if you know the average mass of one unit, a mass \emph{is} a
count. That is the whole trick of the chapter. For atoms the ``average
mass of one unit'' is the tabulated atomic mass: a
\textbf{weighted average} over the natural isotopes,
\[ \bar m = \sum (\text{fractional abundance} \times \text{isotope mass}), \]
measured by a mass spectrometer, where \textbf{peak height gives
abundance}. No single atom has the average mass --- there is no chlorine
atom weighing 35.45 u.

\begin{workedexample}[I do: the atomic mass of silver]
Natural silver is 51.839\% \ce{^{107}Ag} (106.905 u) and 48.161\%
\ce{^{109}Ag} (108.905 u). Compute the tabulated atomic mass.

\[ \bar m = (0.51839)(106.905) + (0.48161)(108.905) \]
\[ \phantom{\bar m} = 55.418 + 52.450 = \mathbf{107.87~u} \]

\textbf{Sanity checks, always both:} the result lies \emph{between} the
two isotope masses, and it leans toward \ce{^{107}Ag} --- the more
abundant isotope. An answer outside 106.905--108.905 means an abundance
was dropped; an answer leaning the wrong way means the abundances were
swapped.
\end{workedexample}

\begin{yourturn}[YOUR TURN 1]
Natural gallium is 60.108\% \ce{^{69}Ga} (68.926 u) and 39.892\%
\ce{^{71}Ga} (70.925 u). Compute the atomic mass, and state both sanity
checks. \workspace[3cm]
\selfcheck{\SI{69.72}{u} --- between the isotope masses, leaning toward
\ce{^{69}Ga}}
\keyonly{\begin{apcanswer}$(0.60108)(68.926) + (0.39892)(70.925)
= 41.430 + 28.294 = \mathbf{69.72}$~u. Between 68.926 and 70.925, and
closer to 68.926 because \ce{^{69}Ga} is the more abundant
isotope.\end{apcanswer}}
\end{yourturn}

% =====================================================================
\section*{\sffamily Ladder 2 \textbullet\ Working backward to an
abundance}
\zsec{3.2}

Given the average and the isotope masses, the abundances follow from one
equation: if $x$ is the fraction of the lighter isotope, the fractions
must sum to 1, so the heavier one is $(1-x)$.

\begin{workedexample}[I do: the isotopes of boron]
Boron has two isotopes, \ce{^{10}B} (10.013 u) and \ce{^{11}B}
(11.009 u), and a tabulated mass of 10.81 u. Find the percent abundances.

\textbf{Set up} with $x$ = fraction of \ce{^{10}B}:
\[ x(10.013) + (1 - x)(11.009) = 10.81 \]

\textbf{Solve:}
\[ 11.009 - 0.996\,x = 10.81 \quad\Longrightarrow\quad
   x = \frac{11.009 - 10.81}{11.009 - 10.013} = \frac{0.199}{0.996}
   = 0.200 \]

So boron is \textbf{20.0\% \ce{^{10}B} and 80.0\% \ce{^{11}B}}.

\textbf{Sanity check:} 10.81 sits much closer to 11.009 than to 10.013,
so the heavy isotope had to dominate --- and it does, four to one.
\end{workedexample}

\begin{yourturn}[YOUR TURN 2]
Rubidium has two isotopes, \ce{^{85}Rb} (84.912 u) and \ce{^{87}Rb}
(86.909 u), and a tabulated mass of 85.47 u. Find both percent
abundances. \workspace[3.4cm]
\selfcheck{72.1\% \ce{^{85}Rb}, 27.9\% \ce{^{87}Rb}}
\keyonly{\begin{apcanswer}$x(84.912) + (1-x)(86.909) = 85.47
\Rightarrow x = (86.909 - 85.47)/(86.909 - 84.912) = 1.439/1.997 =
0.721$. So \textbf{72.1\%} \ce{^{85}Rb} and \textbf{27.9\%} \ce{^{87}Rb}.
Sanity: 85.47 is near 84.912, so the light isotope dominates ---
consistent.\end{apcanswer}}
\end{yourturn}

% =====================================================================
\section*{\sffamily Ladder 3 \textbullet\ The mole}
\zsec{3.3}

One mole is $6.022\times10^{23}$ of anything --- chosen so that one mole
of an element has a mass in grams numerically equal to its atomic mass in
u. Two conversions, and only two:
\[ \text{grams} \xrightarrow{\;\div M\;} \text{moles}
   \xrightarrow{\;\times N_A\;} \text{particles} \]
and each arrow reverses ($\times M$; $\div N_A$) to run the other way.
Avogadro's number converts \textbf{moles to particles}, never grams to
particles directly.

\begin{workedexample}[I do: atoms in a foil strip]
How many aluminum atoms are in a \SI{10.0}{\gram} strip of foil
($M_{\ce{Al}} = \SI{26.98}{\gram\per\mole}$)?

\textbf{Grams to moles first:}
\[ \frac{10.0}{26.98} = \SI{0.371}{\mole}~\ce{Al} \]

\textbf{Then moles to atoms:}
\[ 0.371 \times 6.022\times10^{23} = \mathbf{2.23\times10^{23}~atoms} \]

\textbf{Sense check:} a hand-sized object should hold a substantial
fraction of a mole --- $10^{23}$ is the right neighbourhood. An answer
like $10^{2}$ or $10^{46}$ means a conversion was skipped or doubled.
\end{workedexample}

\begin{yourturn}[YOUR TURN 3]
What is the mass of a sample containing $1.00\times10^{22}$ copper atoms
($M_{\ce{Cu}} = \SI{63.55}{\gram\per\mole}$)? \workspace[2.8cm]
\selfcheck{\SI{1.06}{\gram}}
\keyonly{\begin{apcanswer}$1.00\times10^{22} / 6.022\times10^{23} =
\SI{0.0166}{\mole}$; $0.0166 \times 63.55 = \mathbf{1.06}$~g. The chain
ran particles $\to$ moles $\to$ grams --- the reverse of the worked
example, same two conversions.\end{apcanswer}}
\end{yourturn}

% =====================================================================
\section*{\sffamily Ladder 4 \textbullet\ Molar mass of a compound}
\zsec{3.4}

A compound's molar mass is the sum over the formula, \emph{every
subscript honoured} --- and once you have moles of the compound, the
subscripts also convert to moles of each element inside it.

\begin{workedexample}[I do: aluminum sulfate, down to the atoms]
For \ce{Al2(SO4)3}: (a) the molar mass, (b) the moles in
\SI{10.0}{\gram}, (c) the number of oxygen \emph{atoms} in that sample.

\textbf{(a)} The parentheses multiply everything inside:
\[ M = 2(26.98) + 3(32.07) + 12(16.00)
     = 53.96 + 96.21 + 192.0 = \SI{342.2}{\gram\per\mole} \]
(12 oxygens --- $4 \times 3$ --- not 4. Mis-reading parentheses is the
single commonest molar-mass error.)

\textbf{(b)} $\dfrac{10.0}{342.2} = \SI{0.0292}{\mole}$

\textbf{(c)} Each formula unit carries 12 O atoms:
\[ 0.0292 \times 12 = \SI{0.351}{\mole}~\ce{O}
   \;\Rightarrow\; 0.351 \times 6.022\times10^{23}
   = \mathbf{2.11\times10^{23}~atoms} \]
\end{workedexample}

\begin{yourturn}[YOUR TURN 4]
For \ce{K2Cr2O7} ($M_{\ce{K}} = 39.10$, $M_{\ce{Cr}} = 52.00$):
(a) the molar mass, (b) the moles in \SI{5.00}{\gram}, (c) the number of
oxygen atoms in that sample. \workspace[3.4cm]
\selfcheck{(a) \SI{294.2}{\gram\per\mole} \quad (b) \SI{0.0170}{\mole}
\quad (c) $7.16\times10^{22}$ O atoms}
\keyonly{\begin{apcanswer}(a) $2(39.10) + 2(52.00) + 7(16.00) = 78.20 +
104.0 + 112.0 = \mathbf{294.2}$~g/mol. (b) $5.00/294.2 =
\mathbf{0.0170}$~mol. (c) $0.0170 \times 7 = 0.119$~mol O $\times\,
6.022\times10^{23} = \mathbf{7.16\times10^{22}}$
atoms.\end{apcanswer}}
\end{yourturn}

% =====================================================================
\section*{\sffamily Ladder 5 \textbullet\ Zumdahl's three questions}
\zsec{3.5}

Section 3.5 teaches no chemistry --- it teaches the \emph{approach} that
every multi-step problem in this book expects:

\begin{center}
\fbox{\parbox{0.86\linewidth}{\centering
\textbf{Where am I going?} \quad \textbf{What do I know?} \quad
\textbf{Does my answer make sense?}}}
\end{center}

Plan the route \emph{before} touching the calculator, and check the
result \emph{after}. The middle --- arithmetic --- is the easy part.

\begin{workedexample}[I do: the three questions, worked visibly]
How many hydrogen atoms are in \SI{8.52}{\gram} of ammonia, \ce{NH3}?

\textbf{Where am I going?} From grams of a compound to \emph{atoms of one
element inside it}. Route: grams $\to$ mol \ce{NH3} $\to$ mol H (via the
subscript) $\to$ atoms (via $N_A$). Three hops, planned before any
arithmetic.

\textbf{What do I know?} $M_{\ce{NH3}} = 14.01 + 3(1.008) =
\SI{17.03}{\gram\per\mole}$; each \ce{NH3} carries 3 H.

\textbf{Execute the plan:}
\[ \frac{8.52}{17.03} = \SI{0.500}{\mole}~\ce{NH3}
   \;\xrightarrow{\times 3}\; \SI{1.50}{\mole}~\ce{H}
   \;\xrightarrow{\times N_A}\; \mathbf{9.04\times10^{23}~atoms} \]

\textbf{Does it make sense?} Half a mole of molecules, three H each ---
about 1.5 moles of H, so a bit more than $N_A$ atoms. Yes.
\end{workedexample}

\begin{yourturn}[YOUR TURN 5]
What mass of oxygen is contained in \SI{0.750}{\mole} of \ce{CO2}?
\textbf{First write the route as arrows} (no numbers), then execute, then
state why the size of the answer is reasonable. \workspace[3.2cm]
\selfcheck{route: mol \ce{CO2} $\to$ mol O $\to$ g O; answer
\SI{24.0}{\gram}}
\keyonly{\begin{apcanswer}Route: mol \ce{CO2}
$\xrightarrow{\times 2}$ mol O $\xrightarrow{\times 16.00}$ g O.
$0.750 \times 2 = 1.50$~mol O $\times\, 16.00 = \mathbf{24.0}$~g. Sense:
0.750 mol of \ce{CO2} is about 33 g, and oxygen is $32/44$ of the
compound's mass --- roughly 24 g. Consistent.\end{apcanswer}}
\end{yourturn}

% =====================================================================
\section*{\sffamily Ladder 6 \textbullet\ Percent composition}
\zsec{3.6}

Each element's mass percent is its share of the molar mass:
\[ \%~\text{element} =
   \frac{n \times M_{\text{element}}}{M_{\text{compound}}} \times 100 \]
where $n$ is the subscript. The percents must total 100 (within
rounding) --- that check is free, so always run it.

\begin{workedexample}[I do: aspirin, and a real tablet]
Aspirin is \ce{C9H8O4}. Find the mass percent of each element, then the
mass of carbon in one \SI{500.0}{\milli\gram} tablet.

\textbf{Molar mass:}
$M = 9(12.01) + 8(1.008) + 4(16.00) = 108.09 + 8.064 + 64.00 =
\SI{180.15}{\gram\per\mole}$

\textbf{Each element's share:}
\[ \%\ce{C} = \frac{108.09}{180.15} \times 100 = 60.00\% \qquad
   \%\ce{H} = \frac{8.064}{180.15} \times 100 = 4.48\% \]
\[ \%\ce{O} = \frac{64.00}{180.15} \times 100 = 35.53\% \]

\textbf{Total check:} $60.00 + 4.48 + 35.53 = 100.01$ --- within rounding
of 100 (each percent was rounded to two decimals). A total off by more
than about $\pm 0.05$ is not rounding; it is an arithmetic slip.

\textbf{The tablet:} percent composition is a recipe for \emph{any}
sample size:
\[ m_{\ce{C}} = 0.5000 \times 0.6000 = \SI{0.300}{\gram}~\text{of carbon} \]
\end{workedexample}

\begin{yourturn}[YOUR TURN 6]
Ammonium nitrate fertilizer is \ce{NH4NO3}
($M = \SI{80.05}{\gram\per\mole}$). Find the mass percent of each
element (run the total check), then the mass of nitrogen in a
\SI{25.0}{\gram} scoop. \workspace[4cm]
\selfcheck{35.00\% N, 5.04\% H, 59.96\% O (total 100.00); N in the scoop
\SI{8.75}{\gram}}
\keyonly{\begin{apcanswer}$\%\ce{N} = 2(14.01)/80.05 = 35.00\%$;
$\%\ce{H} = 4(1.008)/80.05 = 5.04\%$;
$\%\ce{O} = 3(16.00)/80.05 = 59.96\%$. Total $100.00$.
$m_{\ce{N}} = 25.0 \times 0.3500 = \mathbf{8.75}$~g. Note both N atoms
count --- the subscripts are read across the whole formula, ammonium and
nitrate alike.\end{apcanswer}}
\end{yourturn}

% =====================================================================
\section*{\sffamily Mastery tracker}

Tick a row only if your YOUR TURN was right on the \emph{first} attempt.
Any unticked row: re-study that ladder's worked example before the next
session, then redo the problem from a blank page.

\renewcommand{\arraystretch}{1.35}
\noindent\begin{tabular}{@{}c p{6.6cm} c p{4.4cm}@{}}
\toprule
\textbf{First try?} & \textbf{Skill} & \textbf{Ladder} &
  \textbf{If not, re-read\ldots}\\
\midrule
$\square$ & weighted-average atomic mass & 1 & silver; both sanity checks \\
$\square$ & abundances from the average & 2 & boron; the $(1-x)$ setup \\
$\square$ & grams--moles--particles & 3 & the two-conversion chain \\
$\square$ & molar mass; atoms inside a sample & 4 & the parentheses rule \\
$\square$ & planning before calculating & 5 & the three questions \\
$\square$ & percent composition & 6 & aspirin; the total check \\
\bottomrule
\end{tabular}

\begin{apcnote}[Scoring yourself honestly]
6/6 first-try: go straight on to the \S3.7--3.11 self-study --- these six
skills are exactly what it assumes. 4--5: solid; redo the missed ladders
tomorrow, not today. 3 or fewer: the gap is usually one of three habits
--- (1) grams never convert to particles directly, only through moles;
(2) subscripts multiply, including through parentheses; (3) every answer
gets a does-it-make-sense pass before you write it down. Fix the habit
and the ladders fall together.
\end{apcnote}

\end{document}
